\documentclass[10pt]{article}
\usepackage[a4paper,margin=25mm]{geometry}
\usepackage{amsmath,amssymb,amsthm,booktabs,array}
\usepackage[T1]{fontenc}
\usepackage{lmodern}
\usepackage[hidelinks]{hyperref}
\newtheorem{theorem}{Theorem}
\newtheorem{lemma}[theorem]{Lemma}
\newtheorem{corollary}[theorem]{Corollary}
\newcommand{\Z}{\mathbb Z}
\newcommand{\dist}{\operatorname{dist}}
\title{Prism-stack shell bounds and proofs of\\Barker's and Spezia's conjectures}
\author{Haobo Ma\\\small ChronoAI Pte Ltd, Singapore
\and Wenlin Zhang\\\small National University of Singapore\\\small ChronoAI Pte Ltd, Singapore}
\date{November 16, 2026}
\begin{document}
\maketitle
\begin{abstract}
Extending Deza and Shtogrin's scale-two metric description
\cite[Tables 3--4]{deza} to exact shell enumeration, we prove
Colin Barker's conjectures in the formula block of OEIS A299289,
revision 16 (formulas dated February 11, 2018), and Stefano Spezia's conjectured exponential
generating function in the same block (June 8, 2024), by expressing distance
in the alternating-prism tiling as layer distance plus distance to an
integer rectangle.
The shell counts are $a(0)=1$ and $2a(n)=13n^2+4$ for positive even $n$,
or $13n^2+3$ for odd $n$.
They imply all of Barker's partial-sum conjectures in A299290,
revision 10, formula block (February 11, 2018).
For arbitrary stacks of square-faced triangular-prism slabs, the metric
gives sharp shell bounds; a nonalternating stack can have the same
coordination sequence at one center but a different sequence at another.
\end{abstract}

\section{The tiling and its graph}
The sequence in A299289~\cite{oeis} counts vertices at a prescribed graph
distance in the uniform tiling \texttt{tsi}. Its initial values were
computed by Davide M. Proserpio with ToposPro. Barker proposed exact
formulas for the shells and for their partial sums~\cite{balls}; Spezia
later proposed an exponential generating function. We obtain these
identities from the sizes of rectangle shells, rather than extending a
finite enumeration.
The source calls Barker's formulas ``Conjectures'' and Spezia's formula
``Conjectured e.g.f.''; the exact lines are identified in the coverage table.
General automata methods establish rationality for periodic graphs
\cite{kopczynski}, whereas the rectangle metric here gives the specific
coefficients directly and also applies to stacks without periodicity.
The metric framework has an earlier source: Deza and Shtogrin
\cite[Table 3, row 27; Table 4, row 32; \S5]{deza} give scale-two
$\ell^1$ embeddings of both the alternating and the general prism stacks.
We give explicit coordinates and use them for shell enumeration and bounds.
Spezia's general OGF--EGF conversion argument~\cite{spezia} explains
the relation between the two series; identifying their coefficients with
the graph-distance shells is the combinatorial step proved here.

The relevant geometry is Gr\"unbaum's tiling 12: slabs of regular
triangular prisms with square boundary faces, with adjacent slab
orientations differing by $90^\circ$~\cite[pp.~52,55]{grunbaum}.
Here is an explicit coordinate model. The graph $G$ has vertices
$(i,j,k)\in\Z^3$. Within each layer $k$, edges are the unit steps in
$i$ or $j$. The upward neighbors of $(i,j,k)$ are
\[
 (i,j,k+1),\qquad
 \begin{cases}(i+1,j,k+1),&k\text{ even},\\
 (i,j+1,k+1),&k\text{ odd}.
 \end{cases}                                                    \tag{1}
\]
All edges are undirected. Put
\[
 \alpha_k=\left\lfloor\frac{k+1}{2}\right\rfloor,\qquad
 \beta_k=\left\lfloor\frac{k}{2}\right\rfloor.
\]
These are floor divisions also for negative $k$.
The Euclidean realization
\[
 (i,j,k)\longmapsto
 \left(i-\frac{\alpha_k}{2},\ j-\frac{\beta_k}{2},\ \frac{\sqrt3}{2}k\right)
                                                                  \tag{2}
\]
identifies this graph with the tiling's edge graph. Indeed, in an even
slab the upper square grid is shifted by half a unit in the first
direction. The triangles between adjacent horizontal rows are
equilateral, and extrusion through a unit interval in the second
direction gives regular triangular prisms. The odd slabs interchange
the two horizontal directions. Each slab is filled face-to-face by
these prisms, its boundary is a square grid, and its edges are exactly
(1) and the horizontal edges. This is the alternating square-faced
prism construction of tiling 12.

Let $a(n)$ be the number of vertices at distance $n$ from the origin,
and let $b(n)=\sum_{r=0}^n a(r)$.

\section{An exact metric and the shell formula}
For integers $c,t$, let $I_c$ be the integer interval with endpoints
$0,c$, and write
\[
 \delta_c(t)=\max(0,\min(0,c)-t)+\max(0,t-\max(0,c)).
\]
Thus $\delta_c(t)$ is the distance from $t$ to $I_c$.

\begin{theorem}\label{main}
For every $(i,j,k)\in\Z^3$,
\[
 \dist_G((0,0,0),(i,j,k))
 =|k|+\delta_{\alpha_k}(i)+\delta_{\beta_k}(j).                 \tag{3}
\]
The graph is vertex-transitive, $a(0)=1$, and at every center
\[
 a(n)=\begin{cases}
 (13n^2+4)/2,&n>0\text{ even},\\
 (13n^2+3)/2,&n\text{ odd}.
 \end{cases}                                                    \tag{4}
\]
\end{theorem}
\begin{proof}
First observe the elementary identity
\[
 H(c,t):=|c|+\delta_c(t)=\max\{|c|,|t|,|t-c|\}.                \tag{5}
\]
It follows by considering $t$ below, inside or above the interval
$I_c$. Since $|\alpha_k|+|\beta_k|=|k|$, the right side $D(i,j,k)$
of (3) equals $H(\alpha_k,i)+H(\beta_k,j)$.
A horizontal edge changes one $t$ by one, so changes $D$ by at most
one. Across an even upward slab, $(\alpha_k,i)$ changes to
$(\alpha_k+1,i+\varepsilon)$, where $\varepsilon=0$ or $1$, while
$(\beta_k,j)$ is fixed. Each of the three absolute values in (5)
changes by at most one. The odd slab has the identical argument in
the second coordinate. Hence $D$ changes by at most one across every
edge. As $D(0,0,0)=0$, every walk to $(i,j,k)$ has length at least $D(i,j,k)$.

For the reverse inequality, take a path monotone in the layer index
from $0$ to $k$. If $k\geq0$, there are $\alpha_k$ first-direction
choices and $\beta_k$ second-direction choices of a shifted upward
edge. Choosing any number of these shifts reaches every pair in
$I_{\alpha_k}\times I_{\beta_k}$ after exactly $k$ steps. If $k<0$,
the downward shifted edges subtract one; there are $|\alpha_k|$ and
$|\beta_k|$ such choices, reaching every pair in the same rectangle
after $|k|$ steps. Select a closest pair to $(i,j)$ in that rectangle
and finish with horizontal edges. This takes exactly $D(i,j,k)$ steps,
proving (3), connectedness, and finiteness of every sphere.

To count a sphere, consider more generally a rectangle
$I_c\times I_e$, with side lengths $u=|c|$, $v=|e|$ in edge units.
The number of integer pairs satisfying
$\delta_c(i)+\delta_e(j)=d$ is
\[
 R_{u,v}(d)=\begin{cases}
 (u+1)(v+1),&d=0,\\
 2(u+v)+4d,&d\geq1.
 \end{cases}                                                    \tag{6}
\]
For $d>0$, the two cases in which one gap is zero contribute
$2(u+1)+2(v+1)$ pairs. Each of the $d-1$ ways to have both gaps
positive contributes four. This also covers zero side lengths.

For $n\geq1$, layers $|k|>n$ contribute nothing. In an interior
layer $|k|<n$, (6) gives
\[
 2|k|+4(n-|k|)=4n-2|k|.
\]
The two extreme layers $k=\pm n$ contribute
$2(\lfloor n/2\rfloor+1)(\lceil n/2\rceil+1)$ in total.
Since $\sum_{|k|<n}|k|=n(n-1)$, summation yields
\[
 a(n)=6n^2-2n+
       2(\lfloor n/2\rfloor+1)(\lceil n/2\rceil+1).             \tag{7}
\]
Substituting $n=2m$ or $n=2m+1$ gives (4). Formula (3) also shows
that the radius-zero sphere contains exactly the origin.

Finally, horizontal integer translations and
$F(i,j,k)=(j,i,k+1)$ are graph automorphisms. In (1), $F$ swaps
the horizontal directions and reverses the parity selecting the shifted
edge; its inverse is $(i,j,k)\mapsto(j,i,k-1)$.
Powers of $F$ reach every layer, and translations reach every vertex
within it. Sphere sizes therefore do not depend on the center.
\end{proof}

\section{Generating functions and partial sums}
\begin{corollary}\label{series}
As formal power series,
\begin{align}
 A(z):=\sum_{n\geq0}a(n)z^n
 &=\frac{1+6z+12z^2+6z^3+z^4}{(1-z)^3(1+z)},                 \tag{8}\\
 E(t):=\sum_{n\geq0}a(n)\frac{t^n}{n!}
 &=\frac{(4+13t+13t^2)\cosh t+(3+13t+13t^2)\sinh t-2}{2}.   \tag{9}
\end{align}
Moreover, $a(n)=2a(n-1)-2a(n-3)+a(n-4)$ for $n\geq5$.
\end{corollary}
\begin{proof}
For every $n\geq0$ we can write
\[
 a(n)=\frac{13}{2}n^2+\frac74+\frac14(-1)^n-\mathbf1_{n=0}.
                                                                  \tag{10}
\]
Using $\sum n^2z^n=z(1+z)/(1-z)^3$, the ordinary series is
\[
 A(z)=\frac{13z(1+z)}{2(1-z)^3}+\frac7{4(1-z)}+
       \frac1{4(1+z)}-1.
\]
Reduction to a common denominator gives (8).
Its denominator is $1-2z+2z^3-z^4$, so comparison of coefficients
above degree four gives the stated recurrence. The initial terms are
$1,8,28,60,106$; the recurrence does not extend to $n=4$, where its
right side would be $105$.
For the exponential series, use
$\sum n^2t^n/n!=(t^2+t)e^t$ in (10). This gives
\[
 E(t)=\frac{(26t^2+26t+7)e^t+e^{-t}-4}{4},
\]
which equals (9) on replacing $e^{\pm t}$ by $\cosh t\pm\sinh t$.
\end{proof}

\begin{corollary}\label{partial}
For every integer $n\geq0$,
\[
 b(n)=\begin{cases}
 (26n^3+39n^2+34n+12)/12,&n\text{ even},\\
 (26n^3+39n^2+34n+9)/12,&n\text{ odd}.
 \end{cases}                                                    \tag{11}
\]
Its ordinary generating function is
\[
 B(z)=\frac{1+6z+12z^2+6z^3+z^4}{(1-z)^4(1+z)},             \tag{12}
\]
and, for $n\geq5$,
\[
 b(n)=3b(n-1)-2b(n-2)-2b(n-3)+3b(n-4)-b(n-5).              \tag{13}
\]
\end{corollary}
\begin{proof}
Summing (10), with
$\sum_{r=0}^n r^2=n(n+1)(2n+1)/6$ and
$\sum_{r=0}^n(-1)^r=(1+(-1)^n)/2$, gives
\[
 b(n)=\frac{13n(n+1)(2n+1)}{12}+
       \frac{7(n+1)}4+\frac{1+(-1)^n}{8}-1.
\]
Expansion proves (11), including $b(0)=1$.
Finite coefficient summation gives $B(z)=A(z)/(1-z)$, proving (12).
The denominator is $1-3z+2z^2+2z^3-3z^4+z^5$; coefficient comparison
for $n\geq5$ proves (13). These are precisely the three formulas
proposed in A299290~\cite{balls}.
\end{proof}

\section{Arbitrary slab orientations}
The same calculation applies when slab directions need not alternate.
Choose any bi-infinite sequence $\varepsilon_k\in\{1,2\}$, and replace
the shifted upward edge in (1) by a shift in coordinate
$\varepsilon_k$. Let $c_0=d_0=0$ and define the integer sequences by
\[
 c_{k+1}-c_k=\mathbf1_{\varepsilon_k=1},\qquad d_k=k-c_k
 \quad(k\in\Z).
\]
This graph is the edge graph of a stack of the same square-faced prism
slabs; it need not be a uniform tiling.

\begin{corollary}\label{stacks}
In this graph the distance to $(i,j,k)$ is
$|k|+\delta_{c_k}(i)+\delta_{d_k}(j)$.
For $n\geq1$, its origin shell size $s(n)$ is
\begin{align}
 s(n)&=6n^2-2n+(|c_n|+1)(|d_n|+1)
                     +(|c_{-n}|+1)(|d_{-n}|+1), \tag{14}\\
 6n^2+2&\leq s(n)\leq a(n). \tag{15}
\end{align}
Both bounds are attained at every radius by fixed choices of the stack.
The same formula holds about any center after shifting the layer origin
and the horizontal coordinates to that center.
An explicit scale-two embedding is
\[
 \Phi(i,j,k)=(k,i,i-c_k,j,j-d_k)\in\Z^5,
 \qquad \|\Phi(p)-\Phi(q)\|_1=2\dist(p,q).
\]
\end{corollary}
\begin{proof}
We have $|c_k|+|d_k|=|k|$. Thus the metric candidate is
$H(c_k,i)+H(d_k,j)$. Each edge changes one of its two summands by
at most one, exactly as in (5). A monotone layer path reaches every
point of $I_{c_k}\times I_{d_k}$; completing horizontally attains the
candidate. The interior-layer count is still $4n-2|k|$, giving
$6n^2-2n$ after summation; the two extreme rectangles give (14).
For nonnegative integers $u,v$ with $u+v=n$,
\[
 n+1\leq(u+1)(v+1)
 \leq(\lfloor n/2\rfloor+1)(\lceil n/2\rceil+1).
\]
Indeed the product is $n+1+uv$, minimized at $uv=0$ and maximized
when $|u-v|\leq1$. Apply this to both extreme rectangles and use (7).
Constant orientations attain the lower bound; alternating orientations
attain the upper. Translating a center to $(0,0,0)$ replaces
$\varepsilon_k$ by the correspondingly shifted sequence, so the
same proof applies there.
For the final identity use
$2H(c,t)=|c|+|t|+|t-c|$ in each coordinate, and note that
$|c_k-c_l|+|d_k-d_l|=|k-l|$. This proves the displayed
embedding for every pair of vertices, recovering the earlier
embeddability statements in~\cite{deza}.
\end{proof}

For a constant orientation the lower bound recovers the known sequence
$h(0)=1$, $h(n)=6n^2+2$ for $n>0$ of the triangular-prism tiling
\texttt{hex} (A005897)~\cite{hex}. Its ball sequence, A005898, is
\[
 H(n)=1+\sum_{r=1}^n(6r^2+2)=n^3+(n+1)^3.
\]
Writing $P(z)=(1+z)(1+4z+z^2)$, their ordinary generating
functions are $P(z)/(1-z)^3$ and $P(z)/(1-z)^4$.
Their exponential generating functions are
$(2+6t+6t^2)e^t-1$ and $(1+8t+9t^2+2t^3)e^t$.
These identities follow by summing $n^2$ or $n^3$ as in
Corollaries~\ref{series}--\ref{partial}; they recover existing formulas,
not further conjectures. Multiplying by the displayed denominators
also recovers the order-three recurrence for $h$ from $n=4$ and the
order-four recurrence for $H$ from $n=4$ recorded in those entries.

The shell sequence does not characterize the alternating stack.
Use the repeating slab word $(1,2,2,1)$, beginning at layer zero.
Every prefix on either side of layer zero contains nearly equal
numbers of the two directions:
each pair of consecutive slabs in those prefixes uses both directions.
Both extreme products in (14) are therefore maximal for every $n$,
and $s(n)=a(n)$. At a center in layer one, however, the two slabs
above have directions $(2,2)$ and those below have directions $(1,1)$.
At radius two (14) gives $24-4+3+3=26$, instead of $a(2)=28$.
This nonuniqueness statement is not a conjecture posed in the source.

\section{Coverage and verification}
The pinned Lean declaration \texttt{result} in~\cite{lean} proves (4)
for the graph (1), at center $(0,0,0)$.
Its supporting lemmas establish the distance and layer-count identities.
The Euclidean identification (2), independence of center, generating
functions, ball formulas, and arbitrary-stack results have written
proofs above; they are not additional conclusions of \texttt{result}.
\begin{center}
\small
\begin{tabular}{@{}p{.48\linewidth}p{.43\linewidth}@{}}
\toprule
Source statement & Coverage here\\
\midrule
A299289 r16, \texttt{\%F} block lines 14--18: Barker's parity formulas & Proved, Theorem~\ref{main}; origin formulas in Lean \texttt{result}\\
A299289 r16, lines 14--18: Barker's ordinary series and recurrence & Proved, Corollary~\ref{series}\\
A299289 r16, line 19: Spezia's exponential series & Proved, Corollary~\ref{series}\\
A299290 r10, lines 11--16: Barker's partial-sum formulas, series and recurrence & Proved, Corollary~\ref{partial}\\
Prism-stack scale-two $\ell^1$ embeddings & Reproved explicitly; prior proof Deza--Shtogrin (1999/2000)\\
Arbitrary-stack shell bounds and nonuniqueness, not posed in OEIS & Proved, Corollary~\ref{stacks} and its example\\
A005897 and A005898: \texttt{hex} shell/ball formulas and series & Prior formulas cited; recovered from Corollary~\ref{stacks}\\
Other cross-referenced tilings besides \texttt{hex} and \texttt{tsi} & Not covered\\
\bottomrule
\end{tabular}
\end{center}
Line numbers refer to the complete plain-text responses linked in
\cite{oeis,balls}, including their introductory headers.
The exact-integer and rational Python checker supplies finite checks
of the graph metric, formulas, series and slab-stack examples; its
README and JSON output specify all ranges. The all-radius arguments
are the proofs above.

\begin{thebibliography}{9}
\raggedright
\bibitem{oeis}
N. J. A. Sloane, \emph{Coordination sequence for ``tsi'' 3D uniform tiling},
The On-Line Encyclopedia of Integer Sequences, A299289, revision 16,
June 8, 2024. Formulas of Colin Barker (February 11, 2018) and
Stefano Spezia (June 8, 2024).
\url{https://oeis.org/A299289}.
\href{https://oeis.org/search?q=id:A299289\&fmt=text}{Plain-text formula block}.
\bibitem{balls}
N. J. A. Sloane, \emph{Partial sums of A299289},
The On-Line Encyclopedia of Integer Sequences, A299290, revision 10,
February 12, 2018. Formulas of Colin Barker (February 11, 2018).
\url{https://oeis.org/A299290}.
\href{https://oeis.org/search?q=id:A299290\&fmt=text}{Plain-text formula block}.
\bibitem{grunbaum}
B. Gr\"unbaum, \emph{Uniform tilings of 3-space},
Geombinatorics \textbf{4} (1994), 49--56; tiling 12 and slab $S_{14}$.
\href{https://faculty.washington.edu/moishe/branko/BG199.Uniform\%20Tilings\%20of\%203-Space.pdf}{Author's article}.
\bibitem{deza}
M. Deza and M. I. Shtogrin, \emph{Uniform partitions of 3-space,
their relatives and embedding}, European Journal of Combinatorics
\textbf{21} (2000), 807--814.
\href{https://doi.org/10.1006/eujc.1999.0385}{doi:10.1006/eujc.1999.0385}.
Preprint arXiv:math/9906034v1, June 6, 1999, Tables 3--4 and \S5.
\url{https://arxiv.org/abs/math/9906034v1}.
\bibitem{kopczynski}
E. Kopczy\'nski, \emph{Coordination Sequences of Periodic Structures
are Rational via Automata Theory}, Acta Crystallographica A
\textbf{78} (2022), 155--157. Theorem 4 and Corollary 5 in the
author's version arXiv:2307.15803v1 (2023).
\href{https://doi.org/10.1107/S2053273322000262}{doi:10.1107/S2053273322000262}.
\url{https://arxiv.org/abs/2307.15803v1}.
\bibitem{hex}
The On-Line Encyclopedia of Integer Sequences, A005897, revision 158
(September 12, 2026), and A005898, revision 140 (October 8, 2026).
The triangular-prism interpretation in A005897 is credited to
N. J. A. Sloane (February 6, 2018); the ordinary generating functions
are credited to Simon Plouffe, and also Colin Barker in A005898.
\url{https://oeis.org/A005897}; \url{https://oeis.org/A005898}.
\bibitem{spezia}
S. Spezia, \emph{Laplace Transforms and the Conversion between Ordinary
and Exponential Generating Functions}, July 28, 2025, updated
September 6, 2025, unnumbered theorem and proof.
\href{https://stefanospezia.altervista.org/laplace-transforms-and-the-conversion-between-ordinary-and-exponential-generating-functions/}{Author's article}.
\bibitem{lean}
\emph{TsiCoordinationSequence.lean}, declaration \texttt{result}.
Fixed source version:
\href{https://github.com/the-omega-institute/trureturing/blob/6974dba01daa1a25cc74b87a86c191e95fe2c681/D5/S3/StatisticalMechanics/TsiCoordinationSequence.lean}{Pinned Lean source}.
\end{thebibliography}
\smallskip
\noindent\small\textcopyright\ 2026 Haobo Ma and Wenlin Zhang.
Licensed under Creative Commons Attribution 4.0 International (CC BY 4.0).
\end{document}
